\documentclass[11pt]{article}
\usepackage[utf8]{inputenc}
\usepackage[T1]{fontenc}
\usepackage[spanish]{babel}
\usepackage[margin=2.3cm]{geometry}
\usepackage{amsmath,amssymb}
\usepackage{enumitem}
\setlength{\parindent}{0pt}
\setlength{\parskip}{6pt}

\begin{document}

\begin{center}
{\Large\bfseries Principios de Macroeconomía}\\[2pt]
{\Large\bfseries Control III}
\end{center}

Miércoles 11/6/2025 \hfill 50 minutos

\textbf{En todas las preguntas sea claro con los gráficos. Nombre los ejes, los puntos de equilibrio y todas las curvas. Diferencia movimientos y desplazamientos de las curvas.}

\begin{enumerate}[leftmargin=*, itemsep=10pt]

\item \textbf{(30 puntos)} Asuma una función de producción Cobb-Douglas $y = f(k) = A k^{0.5}$, donde $A = 1$, además considere los siguientes parámetros que describen la tasa de inversión, el crecimiento poblacional, el crecimiento tecnológico y la depreciación, respectivamente: $s = 0.25$, $n = 0.02$, $\delta = 0{,}08$.

\begin{enumerate}[label=(\alph*), itemsep=3pt]
\item Si la economía cuenta con un capital por trabajador inicial de $k = 4$ en el año 0, ¿cuánto va a ser el producto y la inversión en el año 0?
\item ¿Cuánto va a ser el cambio en el capital, $\dot k$, entre el año 0 y el 1?
\item ¿Y cuánto va a ser el capital por trabajador en el año 1?
\item ¿Cuánto va a ser el producto y la inversión en el año 1?
\item ¿Cuánto va a ser el cambio en el capital, $\dot k$, entre el año 1 y el 2? ¿y cuánto va a ser el capital por trabajador en el año 2?
\item ¿Cuánto va a ser el producto y la inversión en el año 2?
\item ¿Cuánto va a ser el cambio en el capital, $\dot k$, entre el año 2 y el 3? ¿y cuánto va a ser el capital por trabajador en el año 3?
\item ¿Cuánto va a ser el producto y la inversión en el año 3?
\item Calcule el capital y el producto de \textit{Estado Estacionario}, $k^{*}, y^{*}$, en función de los parámetros $s, \delta, n$ y $\alpha$.
\item Reemplace los valores de $s, \delta, n$ y $\alpha$ y calcule el valor de $k^{*}, y^{*}$.
\end{enumerate}

\item \textbf{(10 puntos)} Haga un gráfico (no tiene que ser a escala). Muestre en el gráfico las curvas relevantes y las cantidades de capital, inversión y producto para los años 1, 2, 3 y el Estado Estacionario.

\item \textbf{(10 puntos)} Verdadero o Falso. Según el modelo de Solow visto en clase, en Estado Estacionario las variables per cápita, $y, k$, crecen a tasa 0, mientras que las variables agregadas, $Y, K$, crecen a tasa $(\delta + n)$.

\item \textbf{(10 puntos)} Verdadero o Falso. Según el modelo de Solow visto en clase, un aumento de la tasa de ahorro/inversión, $s$, genera crecimiento en las variables per cápita, $y, k$, en el corto plazo, pero no así en el largo plazo.

\end{enumerate}

\end{document}
